\chapter{Instalación y puesta en marcha}\label{anexo:instalacion}

Ruta comprobada los días 6 y 7 de septiembre de 2026 en dos contenedores nuevos,
sin copiar el entorno Python, el corpus ni las cachés de modelos del equipo.
La evidencia y sus límites están en \href{https://github.com/samubp10/tfg-recomendador-uja/blob/main/docs/experimentos/it54-instalacion.md}{el registro de IT-54}.

\section{Requisitos}

La prueba se hizo en Windows con Docker Desktop 4.69.0, motor Linux 29.4.0,
16 GB de RAM y una NVIDIA RTX 3060 Laptop de 6 GB. Docker disponía de 7,47 GiB
de RAM. Son las condiciones probadas, no unos mínimos medidos. Gemma ocupa
8,1 GB en disco; además hacen falta espacio para imágenes, dependencias,
incrustaciones y datos. Se necesita conexión para instalar, descargar modelos
y extraer la web de la EPSJ.

Esta dotación permitió la consulta web comprobada, pero una prueba de
integración agotó 600 segundos con presión de memoria. Para tandas largas,
comprobar los recursos disponibles y evitar otras consultas simultáneas.
El registro distingue la consulta correcta de la suite con el LLM detenido.

Instalar Docker Desktop con el motor WSL 2 y la compatibilidad NVIDIA descrita
en la \href{https://docs.ollama.com/docker}{documentación de Ollama para Docker}.
Comprobar \texttt{docker version}. Los puertos locales 8000 y 11454 deben estar libres.
Los comandos del equipo anfitrión siguientes son para \textbf{PowerShell}; los de
la sesión del contenedor son para \textbf{sh}. No se necesita Node.js: la interfaz
se sirve desde Python y no tiene un paso de compilación.

\section{Crear el entorno}

En PowerShell, usar nombres nuevos si ya existen estos contenedores:

\begingroup\footnotesize\begin{verbatim}
docker run -d --name it54-app-20260906 `
  -p 127.0.0.1:8000:8001 -p 127.0.0.1:11454:11434 `
  python:3.13-slim sleep infinity
docker run -d --name it54-ollama-20260906 `
  --network container:it54-app-20260906 --gpus all `
  -v it54-modelos-20260906:/root/.ollama `
  ollama/ollama:0.32.14
docker exec it54-ollama-20260906 nvidia-smi
docker exec -it it54-app-20260906 sh
\end{verbatim}\endgroup

Los dos contenedores comparten la red porque la aplicación conecta con Ollama
en \texttt{127.0.0.1:11434}. El volumen conserva los pesos del LLM. Los datos, el índice
y la caché de incrustaciones quedan en el contenedor de la aplicación y sobreviven
a \texttt{docker stop}, pero se perderían al eliminar ese contenedor.

En la sesión \textbf{sh} recién abierta:

\begingroup\footnotesize\begin{verbatim}
apt-get update
apt-get install -y --no-install-recommends git ca-certificates socat
git clone https://github.com/samubp10/tfg-recomendador-uja.git /app
cd /app
git checkout 1ec0484
python -m venv .venv
.venv/bin/python -m pip install torch==2.13.0 \
  --index-url https://download.pytorch.org/whl/cpu
.venv/bin/python -m pip install -e '.[dev,index]' -c constraints.txt
.venv/bin/python -m pip check
\end{verbatim}\endgroup

El commit fija la versión ejecutada en esta comprobación. Para instalar otra
revisión hay que seleccionar su commit y volver a verificarla. El fichero
\texttt{constraints.txt} fija las dependencias; PyTorch se instala primero desde el
repositorio CPU para que las incrustaciones no necesiten CUDA. Ollama utiliza
la GPU por separado. \texttt{pip check} debe terminar sin dependencias rotas.

\section{Configurar las incrustaciones y la base vectorial}

En la misma sesión \textbf{sh}, descargar una vez el modelo elegido:

\begingroup\footnotesize\begin{verbatim}
export TFG_DESCARGAR_MODELO=1
.venv/bin/python - <<'PY'
from tfg_uja.indexacion.incrustaciones import incrustador_de_consultas
print(len(incrustador_de_consultas()(['Ingenieria Informatica'])[0]))
PY
unset TFG_DESCARGAR_MODELO
export HF_HUB_OFFLINE=1
.venv/bin/python - <<'PY'
from tfg_uja.indexacion.incrustaciones import incrustador_de_consultas
print(len(incrustador_de_consultas()(['Ingenieria Informatica'])[0]))
PY
\end{verbatim}\endgroup

Ambas llamadas deben imprimir \texttt{384}. Se usa \texttt{intfloat/multilingual-e5-small},
guardado en \texttt{/root/.cache/huggingface}. La segunda llamada comprueba la carga
desde caché con el Hub desconectado; no prueba que todo el sistema carezca de
acceso a internet. Las variables se describen en la
\href{https://huggingface.co/docs/huggingface_hub/en/package_reference/environment_variables}{documentación de Hugging Face}.
El código añade \texttt{passage: } a los documentos y \texttt{query: } a las preguntas.

Crear el corpus completo y el índice, todavía en \textbf{sh}:

\begingroup\footnotesize\begin{verbatim}
mkdir -p data
.venv/bin/python -m scrapy runspider \
  src/tfg_uja/extraccion/grados_spider.py -O data/grados.json
.venv/bin/python scripts/verificadores/check_dataset.py
.venv/bin/python -m tfg_uja.indexacion.chunker \
  data/grados.json data/chunks.json
.venv/bin/python scripts/verificadores/check_chunks.py
.venv/bin/python scripts/verificadores/check_evalset.py
.venv/bin/python scripts/verificadores/check_guias_pdf.py
.venv/bin/python -m tfg_uja.indexacion.indexer \
  data/chunks.json data/indice_lance
\end{verbatim}\endgroup

El rastreo consulta la web real y respeta sus retardos; no repetirlo para cada
arranque. El verificador de PDF también puede consultar la fuente.
El corpus no está incluido en Git. La web cambia: los recuentos de otra fecha
no tienen por qué coincidir con los del registro de esta prueba.

LanceDB es una biblioteca embebida: no hay servidor, contraseña ni contenedor
de base de datos adicional. El indexador crea \texttt{data/indice\_lance}, con la tabla
\texttt{chunks\_epsj}, y reconstruye la tabla al repetirlo. No cambiar el modelo de
incrustaciones solo en uno de los dos extremos, indexación o consulta.

Salir al anfitrión con \texttt{exit}.

\section{Descargar el LLM y arrancar la aplicación}

En \textbf{PowerShell}:

\begingroup\footnotesize\begin{verbatim}
docker exec it54-ollama-20260906 ollama pull gemma3:12b
docker exec it54-ollama-20260906 ollama list
docker exec -d it54-app-20260906 socat `
  TCP-LISTEN:8001,bind=0.0.0.0,reuseaddr,fork TCP:127.0.0.1:8000
docker exec -it -w /app -e HF_HUB_OFFLINE=1 `
  it54-app-20260906 .venv/bin/python -u -m tfg_uja.aplicacion.servidor
\end{verbatim}\endgroup

\texttt{ollama list} debe mostrar \texttt{gemma3:12b}. El contenedor ya ejecuta el servicio;
no lanzar otro \texttt{ollama serve}. La primera descarga falló por un tiempo de
espera agotado en el alojamiento de los pesos; al recuperar la conexión,
repetir \texttt{ollama pull gemma3:12b} terminó correctamente. No cambiar de modelo
para dar por superada esta comprobación.

Abrir \textbf{http://127.0.0.1:8000}. \texttt{socat} permite llegar al servidor, que escucha
solo en el bucle local del contenedor. El puerto publicado del anfitrión sigue
limitado a \texttt{127.0.0.1}. No publicar el servicio en una interfaz de red externa.
La aplicación usa por defecto \texttt{gemma3:12b}, las incrustaciones anteriores y
\texttt{data/indice\_lance}; no requiere un fichero \texttt{.env} ni claves de servicios remotos.

\section{Comprobar y volver a arrancar}

Preguntar en la interfaz: «¿Cuántos créditos tiene Álgebra en el Grado en
Ingeniería Informática?». Comprobar que llega una respuesta y sus fuentes,
contrastándola con el corpus. Un saludo por sí solo no verifica el LLM.
Para conservar la respuesta HTTP, desde otra ventana de \textbf{PowerShell}:

\begingroup\footnotesize\begin{verbatim}
$cuerpo = @{
  pregunta='¿Cuántos créditos tiene Álgebra en el Grado en Ingeniería Informática?'
} | ConvertTo-Json -Compress
Invoke-WebRequest http://127.0.0.1:8000/api/chat -Method Post `
  -ContentType 'application/json; charset=utf-8' `
  -Body ([System.Text.Encoding]::UTF8.GetBytes($cuerpo)) `
  -TimeoutSec 600 -OutFile consulta.ndjson
\end{verbatim}\endgroup

El fichero contiene los sucesos de la respuesta, uno por línea. La llegada de
HTTP 200 no basta: hay que revisar el contenido y el suceso final \texttt{fin}.

Pruebas del código instalado, sin otra consulta en curso:

\begingroup\footnotesize\begin{verbatim}
docker exec it54-app-20260906 apt-get install -y --no-install-recommends nodejs
docker exec -w /app it54-app-20260906 .venv/bin/python -m pytest `
  --cov=src/tfg_uja --cov=scripts --cov-report=term-missing `
  --cov-fail-under=100
\end{verbatim}\endgroup

Node se instala aquí para las pruebas del cliente JavaScript; no lo usa el
servidor al atender consultas.

Para parar, pulsar Ctrl+C en la sesión del servidor y ejecutar:

\begingroup\footnotesize\begin{verbatim}
docker stop it54-ollama-20260906 it54-app-20260906
\end{verbatim}\endgroup

Para otro uso, arrancar los contenedores:

\begingroup\footnotesize\begin{verbatim}
docker start it54-app-20260906 it54-ollama-20260906
\end{verbatim}\endgroup

Repetir
solo los dos comandos de \texttt{socat} y del servidor del apartado 4. No volver a
clonar, instalar, rastrear o indexar. No ejecutar dos proxies ni dos servidores
a la vez. Para conservar datos fuera del contenedor antes de eliminarlo:

\begingroup\footnotesize\begin{verbatim}
docker cp it54-app-20260906:/app/data ./datos-instalacion
\end{verbatim}\endgroup

Si falta el índice, terminar el apartado 3; si falta la caché de
incrustaciones, repetir su descarga autorizada. Si el LLM no responde,
comprobar la memoria disponible y consultar:

\begingroup\footnotesize\begin{verbatim}
docker logs it54-ollama-20260906
docker exec it54-ollama-20260906 ollama list
\end{verbatim}\endgroup

Un error de conexión del modelo no es una respuesta válida.
